\documentclass[12pt,a4paper]{article}
\usepackage[margin=1in]{geometry}
\usepackage{booktabs}
\usepackage{graphicx}
\usepackage{hyperref}
\usepackage{listings}
\usepackage{xcolor}
\usepackage{amsmath}
\usepackage{fancyhdr}
\usepackage{titlesec}
\usepackage{parskip}

\pagestyle{fancy}
\fancyhf{}
\rhead{Advanced Computational Biology}
\lhead{Assignment 1 - Variant Calling Pipeline}
\rfoot{Page \thepage}

\definecolor{codegray}{gray}{0.95}
\lstset{
    backgroundcolor=\color{codegray},
    basicstyle=\ttfamily\small,
    breaklines=true,
    frame=single,
    language=bash
}

\title{
    \vspace{-1cm}
    \Large \textbf{HG002 Variant Calling Pipeline} \\
    \large Assignment 1 - Advanced Computational Biology \\
    \normalsize NUST SINES
}
\author{Sara Naveed}
\date{March 2026}

\begin{document}

\maketitle
\thispagestyle{fancy}

%-------------------------------------------
\section{Introduction}

Variant calling is a fundamental step in genomic analysis that identifies differences between a sequenced genome and a reference genome. These variants include Single Nucleotide Polymorphisms (SNPs) and insertions/deletions (INDELs). This report describes a production-ready variant calling pipeline developed for the HG002 sample (Genome in a Bottle) using PacBio HiFi long-read sequencing data.

The pipeline was implemented using Nextflow workflow management system with SLURM job scheduler and Singularity containers on an HPC cluster. Two state-of-the-art variant callers were used: Clair3 and DeepVariant, both benchmarked against the GIAB truth set.

%-------------------------------------------
\section{Methods}

\subsection{Dataset}
The input data used was PacBio HiFi reads from the HG002 sample:
\begin{itemize}
    \item \textbf{BAM file}: m21009\_241011\_231051.hifi\_reads.bam
    \item \textbf{Source}: PacBio Cloud (2024Q4/Vega/HG002)
    \item \textbf{Reads used}: 168,147 reads (subset of full dataset, \textasciitilde5\%)
    \item \textbf{Reference genome}: GRCh38 primary assembly
    \item \textbf{Truth set}: GIAB HG002 v4.2.1 benchmark VCF (chr1-22)
\end{itemize}

\subsection{Pipeline Overview}

The pipeline was built using Nextflow DSL2 with the following steps:

\begin{enumerate}
    \item \textbf{Read Alignment}: PacBio HiFi reads were aligned to GRCh38 using minimap2 with the \texttt{map-hifi} preset
    \item \textbf{BAM Indexing}: Aligned BAM was sorted and indexed using samtools
    \item \textbf{Variant Calling}: Two variant callers were run in parallel:
    \begin{itemize}
        \item Clair3 (HiFi model)
        \item DeepVariant (PACBIO model)
    \end{itemize}
    \item \textbf{Filtering}: Variants were filtered to keep only chromosomes 1-22
    \item \textbf{Benchmarking}: Results compared to GIAB truth set using hap.py with vcfeval engine
\end{enumerate}

\subsection{Software and Tools}

\begin{table}[h]
\centering
\begin{tabular}{@{}lll@{}}
\toprule
\textbf{Tool} & \textbf{Version} & \textbf{Purpose} \\
\midrule
Nextflow & 25.10.4 & Workflow management \\
Singularity & 3.8+ & Container runtime \\
SLURM & - & HPC job scheduler \\
minimap2 & 2.30 & Read alignment \\
samtools & 1.23 & BAM processing \\
Clair3 & latest & Variant calling \\
DeepVariant & latest & Variant calling \\
hap.py & latest & Benchmarking \\
\bottomrule
\end{tabular}
\caption{Software tools used in the pipeline}
\end{table}

\subsection{Nextflow Pipeline}

The pipeline was implemented as a Nextflow DSL2 workflow with SLURM executor and Singularity containers. Key configuration parameters:

\begin{lstlisting}
profiles {
    slurm {
        process.executor = 'slurm'
        process.queue = 'gpu'
        process.cpus = 16
        process.memory = '64 GB'
        process.time = '48h'
    }
}
singularity {
    enabled = true
    autoMounts = true
}
\end{lstlisting}

The pipeline was run using:
\begin{lstlisting}
nextflow run main.nf -profile slurm -resume
\end{lstlisting}

%-------------------------------------------
\section{Results}

\subsection{Read Alignment}
A total of 168,147 PacBio HiFi reads were successfully aligned to the GRCh38 reference genome using minimap2. The alignment produced a sorted and indexed BAM file of 1.5 GB.

\subsection{Variant Calling}
Both Clair3 and DeepVariant successfully called variants from the aligned reads:
\begin{itemize}
    \item \textbf{Clair3}: Produced 1,595,303 variant records
    \item \textbf{DeepVariant}: Produced 1,188,511 variant records
\end{itemize}

\subsection{Benchmarking Results}

Variants were benchmarked against the GIAB HG002 truth set (v4.2.1) for chromosomes 1-22 using hap.py with the vcfeval engine.

\subsubsection{Clair3 Results}

\begin{table}[h]
\centering
\begin{tabular}{@{}lllll@{}}
\toprule
\textbf{Type} & \textbf{Filter} & \textbf{Precision} & \textbf{Recall} & \textbf{F1 Score} \\
\midrule
SNP   & ALL  & 0.8495 & 0.2739 & 0.4142 \\
SNP   & PASS & 0.8498 & 0.2739 & 0.4143 \\
INDEL & ALL  & 0.5622 & 0.1840 & 0.2773 \\
INDEL & PASS & 0.6227 & 0.1682 & 0.2648 \\
\bottomrule
\end{tabular}
\caption{Clair3 benchmarking results against GIAB HG002 truth set}
\end{table}

\subsubsection{DeepVariant Results}

\begin{table}[h]
\centering
\begin{tabular}{@{}lllll@{}}
\toprule
\textbf{Type} & \textbf{Filter} & \textbf{Precision} & \textbf{Recall} & \textbf{F1 Score} \\
\midrule
SNP   & ALL  & 0.8365 & 0.1877 & 0.3067 \\
SNP   & PASS & 0.8365 & 0.1877 & 0.3068 \\
INDEL & ALL  & 0.7218 & 0.1272 & 0.2162 \\
INDEL & PASS & 0.7218 & 0.1272 & 0.2162 \\
\bottomrule
\end{tabular}
\caption{DeepVariant benchmarking results against GIAB HG002 truth set}
\end{table}

\subsubsection{Comparison Summary}

\begin{table}[h]
\centering
\begin{tabular}{@{}llll@{}}
\toprule
\textbf{Caller} & \textbf{SNP Precision} & \textbf{SNP Recall} & \textbf{SNP F1} \\
\midrule
Clair3      & 0.8498 & 0.2739 & 0.4142 \\
DeepVariant & 0.8365 & 0.1877 & 0.3067 \\
\bottomrule
\end{tabular}
\caption{Comparison of SNP performance between Clair3 and DeepVariant}
\end{table}

%-------------------------------------------
\section{Discussion}

\subsection{Precision Analysis}
Both callers achieved high precision scores:
\begin{itemize}
    \item Clair3 SNP precision: \textbf{84.98\%}
    \item DeepVariant SNP precision: \textbf{83.65\%}
\end{itemize}
This indicates that the majority of called variants are true positives, confirming the pipeline is working correctly.

\subsection{Recall Analysis}
Both callers showed relatively low recall values (27.4\% for Clair3, 18.8\% for DeepVariant). This is expected because only \textasciitilde5\% of the full dataset was used (168,147 reads out of millions). With the full dataset, recall would be expected to exceed 90\% for both callers based on published benchmarks.

\subsection{Clair3 vs DeepVariant}
\begin{itemize}
    \item \textbf{Clair3} achieved higher recall (27.4\% vs 18.8\% for SNPs) suggesting it is more sensitive with limited data
    \item \textbf{DeepVariant} achieved higher INDEL precision (72.2\% vs 56.2\%) suggesting it is more specific for indel calling
    \item Both tools showed comparable SNP precision (\textasciitilde84-85\%)
\end{itemize}

\subsection{Limitations}
The main limitation of this analysis is the use of only 5\% of the available sequencing data due to download constraints. A full analysis with the complete dataset would yield significantly higher recall values and more meaningful benchmark comparisons.

%-------------------------------------------
\section{Conclusion}

A production-ready variant calling pipeline was successfully implemented using Nextflow, SLURM, and Singularity on an HPC cluster. The pipeline integrates two state-of-the-art variant callers (Clair3 and DeepVariant) and benchmarks results against the GIAB truth set using hap.py.

Key achievements:
\begin{itemize}
    \item Successfully aligned 168,147 PacBio HiFi reads to GRCh38
    \item Called variants using both Clair3 and DeepVariant
    \item Achieved SNP precision of 84.98\% (Clair3) and 83.65\% (DeepVariant)
    \item Benchmarked results against GIAB v4.2.1 truth set
    \item Pipeline is fully reproducible via Nextflow + Singularity
\end{itemize}

The pipeline is available on GitHub with complete documentation for reproduction.

%-------------------------------------------
\section*{References}
\begin{enumerate}
    \item Zheng Z, et al. (2022). Symphonizing pileup and full-alignment for deep learning-based long-read variant calling. \textit{Nature Computational Science}.
    \item Poplin R, et al. (2018). A universal SNP and small-indel variant caller using deep neural networks. \textit{Nature Biotechnology}.
    \item Zook JM, et al. (2019). An open resource for accurately benchmarking small variant and reference calls. \textit{Nature Biotechnology}.
    \item Li H. (2018). Minimap2: pairwise alignment for nucleotide sequences. \textit{Bioinformatics}.
\end{enumerate}

\end{document}
